\chapter*{Lời nói đầu}
\addcontentsline{toc}{chapter}{Lời nói đầu}

Trong chương trình đào tạo ngành Công nghệ Thông tin, học phần \textbf{Cấu trúc dữ liệu và Giải thuật} đóng vai trò là một trong những môn học nền tảng cốt lõi nhất. Môn học không chỉ trang bị tư duy logic trừu tượng, phương pháp phân tích độ phức tạp thời gian và không gian mà còn rèn luyện kỹ năng hiện thực hóa các cấu trúc lưu trữ dữ liệu phi tuyến tính trong bộ nhớ máy tính.

Đề tài \textbf{"Quản lý danh mục sách thư viện bằng cây nhị phân tìm kiếm (BST)"} được nhóm 04 lựa chọn nhằm áp dụng lý thuyết cây nhị phân vào một bài toán thực tiễn quen thuộc trong công tác quản lý biên mục thư viện. Qua đề tài này, nhóm hướng tới các mục tiêu:
\begin{enumerate}
    \item Hiện thực hóa hoàn chỉnh cấu trúc cây nhị phân tìm kiếm động bằng ngôn ngữ C++17, đảm bảo an toàn bộ nhớ (Memory Safety) và không để xảy ra rò rỉ bộ nhớ.
    \item Vận dụng các thuật toán cốt lõi: Chèn node, Tìm kiếm nhị phân $O(h)$, Xóa node triệt để qua 3 trường hợp và duyệt cây theo các thứ tự NLR, LNR, LRN.
    \item Xây dựng ứng dụng hoàn chỉnh với tính năng đọc/ghi file CSV bền vững và giao diện dòng lệnh thân thiện.
    \item Rèn luyện kỹ năng làm việc nhóm hiện đại, quản lý mã nguồn bằng Git, quản lý tác vụ bằng Jira và trình bày báo cáo học thuật chuyên nghiệp bằng \LaTeX.
\end{enumerate}

Nhóm 04 xin gửi lời cảm ơn chân thành tới Quý Thầy/Cô bộ môn đã tận tình hướng dẫn, truyền đạt kiến thức và định hướng phương pháp tư duy giải thuật trong suốt quá trình học tập. Dù đã rất nỗ lực, bài báo cáo khó tránh khỏi những thiếu sót, nhóm rất mong nhận được những ý kiến đóng góp quý báu từ Quý Thầy/Cô.

\begin{flushright}
    \textit{Hà Nội / TP. Hồ Chí Minh, ngày 25 tháng 10 năm 2026}\\
    \textbf{Thay mặt nhóm tác giả}\\
    \textbf{Trưởng nhóm: Trần Ninh Thuận}
\end{flushright}
